\documentclass[10pt,letterpaper]{article}
\usepackage[T1]{fontenc}
\usepackage{amssymb}
\usepackage{amsmath}
\usepackage{braket}
\usepackage{enumerate}
\usepackage[margin=0.5in]{geometry}
\title{Introduction to Quantum Electronic Devices and Applications Qubit Technology Research Paper Proposal}
\author{Jeremy Chen}
\date{03/10/2026}
\begin{document}
	\maketitle

\section*{Proposal}
The modality of qubits that I am selecting is neutral atoms, a method of quantum computing that relies on optical tweezers to trap massive arrays of Rydberg atoms that then act as qubits. Optical tweezers are lasers that use the radiation pressure of photons to trap atomic scale particles. In the case of a neutral atom quantum computer, they are Rydberg atoms, an atom with electrons that are in a very high excited state. Rydberg atoms are more sensitive to electromagnetic fields, hence their utility in being tweezed. 

The physical model of a neutral atom quantum computer is represented by the following Hamiltonian
$$\hat{H} = \hat{H}_{1} + \hat{H}_{2} + V_{rr}(\ket{r}_{1}\bra{r} \otimes \ket{r}_{2}\bra{r})$$
where $\hat{H}_{i}$ is the Hamiltonian of an individual qubit and $V_{rr}$ is the entanglement between the two qubits. This entanglement process is physically achieved through Rydber interactions, which rely on the large electric dipole moments of Rydberg atoms. Perturbation theory shows that this large electric dipole creates strong interaction between two adjacent Rydberg atoms. 

Single qubit gates, on the other hand, are achieved with applying a microwave pulse at the Rydberg atom, which generates rotation much like superconducting qubits. Another method of generating a gate is through laser pulses. 

\section*{References}
\begin{enumerate}[1. ]
	
\item Manetsch, H.J., Nomura, G., Bataille, E. et al. A tweezer array with 6,100 highly coherent atomic qubits. Nature 647, 60–67 (2025). https://doi.org/10.1038/s41586-025-09641-4

\item { \tt arXiv:0909.4777 [quant-ph] }

\item { \tt arXiv:quant-ph/9806021 }

\end{enumerate}

\end{document}